\section{Related Work}\label{sec:related}

Recent advances in artificial intelligence and image-based analysis have
substantially expanded the diagnostic capabilities of oral health systems.
CNNs and, more recently, Transformer-based architectures have demonstrated
strong performance on a range of oral lesion classification and cancer
detection tasks. The existing literature, however, is heavily skewed toward two
narrow problem formulations: binary malignant-versus-benign discrimination, and
detection of a single dominant disease category. Multi-class classification of
clinically heterogeneous oral conditions remains comparatively under-explored,
despite its closer alignment with real screening practice. Compounding this,
most reported architectures inherit attention modules and hyper-parameters from
generic benchmarks without systematic ablation, and few studies pair their
quantitative results with explainability evidence that the model attends to
clinically relevant tissue. Object-detection frameworks have been adapted for
lesion localization but produce only bounding boxes and do not characterize
sub-class identity, while a separate line of work targets pixel-level
segmentation. Both directions are valuable but address questions
orthogonal to the multi-class diagnostic categorization that the present study
targets. The following subsections review each strand in turn and then
position the present work against the gaps that emerge.

\subsection{Binary Oral Cancer Detection}

A substantial body of work has approached oral cancer as a binary screening
problem. The CNN-based framework reported by Jubair \emph{et al.}~\cite{ref5}
used transfer learning to compensate for limited clinical data and reported
approximately $85\,\%$ classification accuracy (with a $95\,\%$ confidence
interval of $81$--$90\,\%$) between malignant and non-malignant lesions,
demonstrating discriminative capacity sufficient for screening. The binary
formulation, however, precludes sub-type differentiation, and the model does
not characterize lesion morphology.

Subsequent efforts have prioritized computational efficiency. The SE-MobileViT
architecture~\cite{ref16} integrates Squeeze-and-Excitation modules into a
MobileViT backbone for channel-wise attention, achieving $98.39\,\%$ binary
accuracy with reduced inference cost suitable for mobile deployment. The same
trade-off applies: the model establishes the presence or absence of malignancy
but does not differentiate among clinically distinct oral conditions.

A complementary direction has combined hand-crafted and learned
representations. The fusion of Local Binary Pattern (LBP) descriptors with deep
CNN features has also been explored for oral cancer detection~\cite{ref17}. The
contribution remains at the representation level --- the model still maps each
image to a binary outcome.

Across this family of methods, the dominant limitation is not accuracy but
scope. Binary screening is insufficient when downstream clinical decisions
depend on identifying which oral condition is present, since ulcerative,
inflammatory, and neoplastic lesions warrant distinct management pathways.

\subsection{Recurrent Aphthous Ulcer-Focused Models}

A smaller body of work has targeted recurrent aphthous ulcer (RAU)
specifically. The CNN-based framework reported in~\cite{ref3} used transfer
learning with ResNet variants to distinguish ulcer and non-ulcer classes on a
modest clinical dataset, achieving above $85\,\%$ accuracy. While the result
demonstrates feasibility under data-scarce conditions, the formulation is
binary and the small dataset constrains generalization to broader oral disease
populations.

A recent systematic review by Tiwari \emph{et al.}~\cite{ref23} surveyed $23$
artificial-intelligence studies on the diagnosis of mouth ulcers --- spanning
oral lichen planus (OLP), recurrent aphthous stomatitis, leukoplakia, and
related conditions --- and reported accuracies ranging from $71\,\%$ to
$100\,\%$ depending on dataset size and architecture choice. The review
highlights two persistent gaps: limited class balance across datasets, and the
near-absence of explainability evidence. These gaps directly motivate the dual
emphasis of the present study on multi-class differentiation and visual
interpretability.

\subsection{Semantic Segmentation of Oral Lesions}

A complementary line of research has approached oral lesions as a pixel-level
or instance-level segmentation problem. The CLASEG framework~\cite{ref15}
addresses differential diagnosis of oral lesions across $14$ lesion classes by
coupling an EfficientNet-B3 classifier with a ResNet-101-based Mask R-CNN
instance-segmentation branch; it reports a classification accuracy of
$74.49\,\%$ and a segmentation $\mathrm{AP}_{50}$ of $72.18$. The work
demonstrates that joint classification and lesion delineation is feasible over
a broad lesion taxonomy, although the moderate classification accuracy
indicates that fine-grained discrimination among many visually similar classes
remains difficult.

A more recent hybrid CNN-Transformer design~\cite{ref26} integrates high-order
focus convolution with edge-aware modules and introduces Sobel-based edge
enhancement to improve boundary delineation. The proposed model, HF-UNet, is
trained on a curated oral-ulcer dataset and reaches a Dice score near $0.80$,
with limited dataset size identified as the principal constraint on further
gains.

Segmentation and classification address different clinical questions:
segmentation localizes lesion extent at pixel or instance granularity, whereas
classification identifies the disease class. The present work focuses on the
latter, where multi-class differentiation among visually similar conditions
remains the dominant challenge.

\subsection{Multi-Class Oral Disease Classification}

The closest prior work to the present study is multi-class oral disease
classification. Rashid \emph{et al.}~\cite{ref19} reported a strong single-model
result, classifying a seven-class oral disease dataset of intraoral images with
an InceptionResNetV2 network under a unified preprocessing protocol, reporting
approximately $99.5\,\%$ accuracy and emphasizing generalization across patient
demographics. While the result indicates that high-accuracy multi-class oral
disease classification is attainable, the study treats backbone selection as
the principal design axis and does not investigate how attention mechanisms
should be integrated within the backbone. We note that this single-model
literature result reports a subtype accuracy comparable to the figure obtained
here on a related oral disease collection; accordingly, every claim of a
``highest subtype accuracy'' in this paper is scoped to the controlled
nine-baseline benchmark trained from scratch under one matched recipe, and is
never advanced as a record against the broader literature.

A related ensemble line of work, reported as MODC-SET~\cite{ref22}, combines
MobileNetV2, InceptionResNetV2, and ResNet50 with an XGBoost meta-classifier and
is evaluated on a separately curated collection of seven oral disease
categories, where it is reported to reach roughly $99\,\%$ overall accuracy with
feature fusion aiding discrimination between visually similar lesions. The
architectural choices --- which backbones to ensemble, and why XGBoost as the
meta-classifier --- are not subjected to ablation, and the substantial
cumulative parameter count of the combined model raises practical deployment
concerns. Because that work is evaluated on a different seven-class collection,
it does not constitute a controlled comparator for the present study; we cite
it only to situate the magnitude of multi-class results in the literature. The
present study addresses the orthogonal question of how an efficient
single-model alternative can be designed in a principled, ablation-driven way.

\subsection{Transformer and Hybrid Architectures}

Vision Transformer (ViT) architectures have also been explored for oral lesion
analysis. Chilet-Martos \emph{et al.}~\cite{ref18} compared a ViT-based pipeline
--- combining DETR, the Segment Anything Model, and a ViT classifier ---
against radiomics-based baselines, with the ViT pipeline reaching an accuracy of
approximately $0.93$; the study is framed primarily as a model-type comparison
between deep representation learning and hand-crafted radiomic features.
Hybrid CNN-Transformer designs that couple a DeiT branch with a CoAtNet branch
have also been explored for image-based disease classification~\cite{ref21},
typically taking the hybrid attention configuration as given rather than testing
it through ablation. A distinct line of work has explored multimodal large
language models as diagnostic assistants --- the evaluation in~\cite{ref24}
tested ChatGPT-5's ability to differentiate OLP, oral lichenoid lesions, and
squamous-cell carcinoma developing over lichen planus on biopsy-confirmed cases,
finding competitive but sub-expert performance that lagged human experts at the
Top-1 differential. That framework emphasizes diagnostic reasoning over spatial
precision and complements rather than replaces dedicated vision models.

\subsection{Position of the Present Work}

The literature reviewed above reveals three gaps that motivate the present
study. \emph{First}, binary formulations dominate; the few multi-class systems
that exist~\cite{ref19,ref22} do not subject their architectural decisions to
ablation. \emph{Second}, attention modules are stacked without principled
justification --- the literature offers no empirically grounded rule for how
attention components should be combined on this distribution. \emph{Third},
explainability evidence in the multi-class oral disease setting is sparse,
leaving the field's high-accuracy claims clinically unverified.

The present work addresses all three gaps. We propose a parameter-efficient
backbone with a novel AttentionHub module whose design is derived from a
seven-cell ablation rather than imported from a generic benchmark; we benchmark
it against nine standard CNN and Transformer baselines under a single matched
training recipe with all models trained from scratch; and we accompany every
model with Grad-CAM++ and LIME panels for both the binary and the seven-class
heads.
